% GENERATED FILE. Edit canonical lessons, then run npm run generate:books.
% canonical-source-hash: fnv1a64:da73a50506652216
% canonical-lessons: TA-C482-culukku, TA-C482-picaku, TA-C482-atipatu, TA-C482-muri, TA-C482-kaci

\chapter{Sprains and Breaks}
\label{ch:ta-sprains-and-breaks}
\hlchaptermodality{482}

\begin{chapteropening}
\textbf{By the end of this chapter:} \emph{I can talk about sprains, dislocations, getting hurt and broken bones.}

Complete the last lesson of chapter 482: I can talk about sprains, dislocations, getting hurt and broken bones.
\end{chapteropening}

\section[\texorpdfstring{\ta{சுளுக்கு} (cuḷukku)}{சுளுக்கு (cuḷukku)}]{\ta{சுளுக்கு} (cuḷukku) — a sprain}

\label{lesson:TA-C482-culukku}



\begin{quote}
Before the new one: say the Tamil for to trip, then the Tamil for to falter.
\end{quote}

\subsection*{You'll want to know: \ta{சுளுக்கு}}
\textbf{\ta{சுளுக்கு}} — \emph{cuḷukku} — \textquotedblleft{}a sprain\textquotedblright{}.

\begin{quote}
\ta{விளையாடும்போது} \ta{எனக்குக்} \ta{காலில்} \ta{சுளுக்கு} \ta{ஏற்பட்டது}.
\end{quote}

\begin{quote}
\emph{viḷaiyāṭumpōtu eṉakkuk kālil cuḷukku ēṟpaṭṭatu.}
\end{quote}

\begin{quote}
I sprained my foot while playing.
\end{quote}

Also a verb: \ta{கால்} \ta{சுளுக்கியது}, the foot got sprained.

\begin{cousinweb}
\textbf{\ta{சுளுக்கு}} (\emph{cuḷukku}) is native Dravidian (DEDR 2702); the entry also has \ta{உளுக்கு} \emph{uḷukku} for the same.
\end{cousinweb}

\begin{grammarlens}[title={What you've built}]
One more word for describing what you lived and felt.
\end{grammarlens}

\subsection*{Your turn}
\begin{itemize}
\raggedright
  \item \emph{Say it:} \emph{cuḷukku}
  \item \emph{Say it:} \emph{cuḷukku}, once more
  \item \emph{Say it:} \emph{viḷaiyāṭumpōtu eṉakkuk kālil cuḷukku ēṟpaṭṭatu.} — I sprained my foot while playing.
  \item \emph{Recall:} say the Tamil for to trip, then the Tamil for to falter, then say \emph{cuḷukku} again
\end{itemize}

\begin{tcolorbox}[breakable,colback=teal!4,colframe=teal!35!black,title={Before you move on}]
What does \ta{சுளுக்கு} mean? (\textquotedblleft{}A sprain\textquotedblright{}.) Say it once more.
\end{tcolorbox}
\section[\texorpdfstring{\ta{பிசகு} (picaku)}{பிசகு (picaku)}]{\ta{பிசகு} (picaku) — to be dislocated; to go wrong}

\label{lesson:TA-C482-picaku}



\begin{quote}
Before the new one: say the Tamil for to falter, then the Tamil for a sprain.
\end{quote}

\subsection*{You'll want to know: \ta{பிசகு}}
\textbf{\ta{பிசகு}} — \emph{picaku} — \textquotedblleft{}to be dislocated; to go wrong\textquotedblright{}.

\begin{quote}
\ta{விழுந்ததில்} \ta{அவருடைய} \ta{தோள்} \ta{பிசகியது}.
\end{quote}

\begin{quote}
\emph{viḻuntatil avaruṭaiya tōḷ picakiyatu.}
\end{quote}

\begin{quote}
His shoulder was dislocated in the fall.
\end{quote}

Past (it …): \ta{பிசகியது}, \emph{picakiyatu}.

Said of the joint, so the past is \ta{பிசகியது}. Of plans too: \ta{கணக்கு} \ta{பிசகியது}, the sums went wrong.

\begin{cousinweb}
\textbf{\ta{பிசகு}} (\emph{picaku}) is native Dravidian (DEDR 4131), from a false step to a slipped joint.
\end{cousinweb}

\begin{grammarlens}[title={What you've built}]
One more word for describing what you lived and felt.
\end{grammarlens}

\subsection*{Your turn}
\begin{itemize}
\raggedright
  \item \emph{Say it:} \emph{picaku}
  \item \emph{Say it:} \emph{picaku}, once more
  \item \emph{Say it:} \emph{viḻuntatil avaruṭaiya tōḷ picakiyatu.} — His shoulder was dislocated in the fall.
  \item \emph{Recall:} say the Tamil for to falter, then the Tamil for a sprain, then say \emph{picaku} again
\end{itemize}

\begin{tcolorbox}[breakable,colback=teal!4,colframe=teal!35!black,title={Before you move on}]
What does \ta{பிசகு} mean? (\textquotedblleft{}To be dislocated; to go wrong\textquotedblright{}.) Say it once more.
\end{tcolorbox}
\section[\texorpdfstring{\ta{அடிபடு} (aṭipaṭu)}{அடிபடு (aṭipaṭu)}]{\ta{அடிபடு} (aṭipaṭu) — to get hurt}

\label{lesson:TA-C482-atipatu}



\begin{quote}
Before the new one: say the Tamil for a sprain, then the Tamil for to be dislocated.
\end{quote}

\subsection*{You'll want to know: \ta{அடிபடு}}
\textbf{\ta{அடிபடு}} — \emph{aṭipaṭu} — \textquotedblleft{}to get hurt\textquotedblright{}.

\begin{quote}
\ta{விளையாடும்போது} \ta{அவனுக்குத்} \ta{தலையில்} \ta{அடிபட்டது}.
\end{quote}

\begin{quote}
\emph{viḷaiyāṭumpōtu avaṉukkut talaiyil aṭipaṭṭatu.}
\end{quote}

\begin{quote}
He hurt his head while playing.
\end{quote}

Past (I …): \ta{அடிபட்டேன்}, \emph{aṭipaṭṭēṉ}.

Often with the dative: \ta{அவனுக்கு} \ta{அடிபட்டது}, he got hurt.

\begin{cousinweb}
\textbf{\ta{அடிபடு}} (\emph{aṭipaṭu}) is built from native words: \ta{அடி} \emph{aṭi}, \textquotedblleft{}to beat; a blow\textquotedblright{} (DEDR 77), and \ta{படு}, \textquotedblleft{}to undergo\textquotedblright{}.
\end{cousinweb}

\begin{grammarlens}[title={What you've built}]
One more word for describing what you lived and felt.
\end{grammarlens}

\subsection*{Your turn}
\begin{itemize}
\raggedright
  \item \emph{Say it:} \emph{aṭipaṭu}
  \item \emph{Say it:} \emph{aṭipaṭu}, once more
  \item \emph{Say it:} \emph{viḷaiyāṭumpōtu avaṉukkut talaiyil aṭipaṭṭatu.} — He hurt his head while playing.
  \item \emph{Recall:} say the Tamil for a sprain, then the Tamil for to be dislocated, then say \emph{aṭipaṭu} again
\end{itemize}

\begin{tcolorbox}[breakable,colback=teal!4,colframe=teal!35!black,title={Before you move on}]
What does \ta{அடிபடு} mean? (\textquotedblleft{}To get hurt\textquotedblright{}.) Say it once more.
\end{tcolorbox}
\section[\texorpdfstring{\ta{முறி} (muṟi)}{முறி (muṟi)}]{\ta{முறி} (muṟi) — to snap; to break}

\label{lesson:TA-C482-muri}



\begin{quote}
Before the new one: say the Tamil for to be dislocated, then the Tamil for to get hurt.
\end{quote}

\subsection*{You'll want to know: \ta{முறி}}
\textbf{\ta{முறி}} — \emph{muṟi} — \textquotedblleft{}to snap; to break\textquotedblright{}.

\begin{quote}
\ta{மரத்திலிருந்து} \ta{விழுந்து} \ta{அவனுடைய} \ta{கை} \ta{முறிந்தது}.
\end{quote}

\begin{quote}
\emph{marattiliruntu viḻuntu avaṉuṭaiya kai muṟintatu.}
\end{quote}

\begin{quote}
He fell from a tree and broke his arm.
\end{quote}

Past (it …): \ta{முறிந்தது}, \emph{muṟintatu}.

Of bones and branches: \ta{கிளை} \ta{முறிந்தது}, the branch snapped.

\begin{cousinweb}
\textbf{\ta{முறி}} (\emph{muṟi}) is native Dravidian: DEDR 4975 lists the verb, spelled \emph{muri}, \textquotedblleft{}to break off, to snap\textquotedblright{}.
\end{cousinweb}

\begin{grammarlens}[title={What you've built}]
One more word for describing what you lived and felt.
\end{grammarlens}

\subsection*{Your turn}
\begin{itemize}
\raggedright
  \item \emph{Say it:} \emph{muṟi}
  \item \emph{Say it:} \emph{muṟi}, once more
  \item \emph{Say it:} \emph{marattiliruntu viḻuntu avaṉuṭaiya kai muṟintatu.} — He fell from a tree and broke his arm.
  \item \emph{Recall:} say the Tamil for to be dislocated, then the Tamil for to get hurt, then say \emph{muṟi} again
\end{itemize}

\begin{tcolorbox}[breakable,colback=teal!4,colframe=teal!35!black,title={Before you move on}]
What does \ta{முறி} mean? (\textquotedblleft{}To snap; to break\textquotedblright{}.) Say it once more.
\end{tcolorbox}
\section[\texorpdfstring{\ta{கசி} (kaci)}{கசி (kaci)}]{\ta{கசி} (kaci) — to ooze; to seep}

\label{lesson:TA-C482-kaci}



\begin{quote}
Before the new one: say the Tamil for to get hurt, then the Tamil for to snap.
\end{quote}

\subsection*{You'll want to know: \ta{கசி}}
\textbf{\ta{கசி}} — \emph{kaci} — \textquotedblleft{}to ooze; to seep\textquotedblright{}.

\begin{quote}
\ta{காயத்திலிருந்து} \ta{இரத்தம்} \ta{கசிந்தது}.
\end{quote}

\begin{quote}
\emph{kāyattiliruntu irattam kacintatu.}
\end{quote}

\begin{quote}
Blood oozed from the wound.
\end{quote}

Past (it …): \ta{கசிந்தது}, \emph{kacintatu}.

Also of damp walls and of tears.

\begin{cousinweb}
\textbf{\ta{கசி}} (\emph{kaci}) is native Dravidian (DEDR 1091), of moisture seeping through a wall or eyes welling up.
\end{cousinweb}

\begin{grammarlens}[title={What you've built}]
One more word for describing what you lived and felt.
\end{grammarlens}

\subsection*{Your turn}
\begin{itemize}
\raggedright
  \item \emph{Say it:} \emph{kaci}
  \item \emph{Say it:} \emph{kaci}, once more
  \item \emph{Say it:} \emph{kāyattiliruntu irattam kacintatu.} — Blood oozed from the wound.
  \item \emph{Recall:} say the Tamil for to get hurt, then the Tamil for to snap, then say \emph{kaci} again
\end{itemize}

\begin{tcolorbox}[breakable,colback=teal!4,colframe=teal!35!black,title={Before you move on}]
What does \ta{கசி} mean? (\textquotedblleft{}To ooze; to seep\textquotedblright{}.) Say it once more.
\end{tcolorbox}
